\documentclass[11pt]{article}

\usepackage[margin=1in]{geometry}
\usepackage[T1]{fontenc}
\usepackage[utf8]{inputenc}
\usepackage{lmodern}
\usepackage{amsmath, amssymb}
\usepackage{booktabs}
\usepackage{longtable}
\usepackage{array}
\usepackage{hyperref}

\hypersetup{
    colorlinks=true,
    linkcolor=blue,
    urlcolor=blue,
    pdftitle={CCP Margin Engine Explainer}
}

\title{CCP Margin Adequacy, Backtesting, and Risk Governance Engine\\
\large Detailed LaTeX Explanation with Mathematical Concepts}
\author{}
\date{\today}

\begin{document}

\maketitle
\tableofcontents
\newpage

\section{Executive Summary}

The CCP Margin Adequacy, Backtesting, and Risk Governance Engine is a stylised risk platform that evaluates whether a clearing member's posted collateral is sufficient to absorb liquidation-adjusted losses under both ordinary and stressed market conditions. The engine is deliberately broader than a stand-alone Value-at-Risk calculation. It combines:

\begin{itemize}
    \item synthetic market and portfolio generation,
    \item full revaluation pricing for futures and listed options,
    \item historical simulation and stressed Value-at-Risk,
    \item liquidity and concentration add-ons,
    \item adequacy controls and margin-call logic,
    \item backtesting and data-quality controls,
    \item rule-based governance escalations, and
    \item committee-ready reporting.
\end{itemize}

The core business question is:
\[
\text{Is posted collateral sufficient to cover liquidation-adjusted loss for each member on each day?}
\]

The engine is educational rather than production-grade, but it is organised in a production-style architecture. The quantitative point is not merely to estimate market risk, but to transform risk estimates into operational and governance decisions.

\section{Notation and System Structure}

\subsection{Indices and State Variables}

Let:
\begin{align*}
m &\in \{1,\dots,M\} && \text{index clearing members},\\
i &\in \{1,\dots,I\} && \text{index instruments},\\
j &\in \{1,\dots,J\} && \text{index risk factors},\\
t &\in \{1,\dots,T\} && \text{index business dates}.
\end{align*}

The key state variables are:
\begin{align*}
S_{j,t} &:= \text{spot level of risk factor } j \text{ at date } t,\\
r_{j,t} &:= \log(S_{j,t}/S_{j,t-1}) \text{ one-day log return},\\
q_{m,i,t} &:= \text{position quantity of member } m \text{ in instrument } i,\\
V_{m,t} &:= \text{portfolio value of member } m \text{ at date } t,\\
\Pi_{m,t} &:= \text{one-day profit and loss of member } m,\\
\text{IM}_{m,t} &:= \text{required initial margin},\\
\text{PM}_{m,t} &:= \text{posted collateral after haircut}.
\end{align*}

\subsection{Architecture}

The implementation is separated into the following modules.

\begin{center}
\begin{tabular}{>{\raggedright\arraybackslash}p{0.22\textwidth} >{\raggedright\arraybackslash}p{0.46\textwidth} >{\raggedright\arraybackslash}p{0.22\textwidth}}
\toprule
Module & Function & Main Files \\
\midrule
Data and portfolio & Generate synthetic market data, instruments, member types, positions, collateral & \texttt{src/portfolio.py}, \texttt{src/data\_loader.py} \\
Pricing and P\&L & Full revaluation of futures and options; daily P\&L generation & \texttt{src/pricing.py} \\
Scenario engine & Historical and stressed return scenario construction & \texttt{src/scenarios.py} \\
Margin methodology & Historical VaR, stressed VaR, liquidity and concentration add-ons & \texttt{src/margin.py}, \texttt{src/liquidity.py}, \texttt{src/concentration.py} \\
Adequacy controls & Coverage ratios, traffic lights, margin call logic & \texttt{src/controls.py} \\
Monitoring & Backtesting exceptions and data-quality checks & \texttt{src/backtesting.py} \\
Governance & Escalation rules and breach handling & \texttt{src/escalation.py} \\
Reporting & Daily, weekly, and committee-style outputs & \texttt{src/reporting.py} \\
\bottomrule
\end{tabular}
\end{center}

\section{Synthetic Data Layer}

\subsection{Risk Factors}

The platform models eight primary risk factors:
\[
\{\text{SPX}, \text{NDX}, \text{RTY}, \text{TY}, \text{FV}, \text{US}, \text{VIX}, \text{CL}\}.
\]

These cover equity-index, rates, volatility, and commodity risk. The synthetic market-data generator uses fat-tailed shocks rather than Gaussian shocks. Specifically, each factor is generated using Student-\(t\) innovations with \(\nu=5\) degrees of freedom:
\[
\varepsilon_{j,t} \sim t_{\nu}, \qquad \nu=5.
\]

If \(\sigma_j^{\text{ann}}\) denotes the annual volatility assigned to factor \(j\), the daily-scale volatility is:
\[
\sigma_j^{\text{day}} = \frac{\sigma_j^{\text{ann}}}{\sqrt{252}}.
\]

The synthetic one-day return is then:
\[
r_{j,t} = \sigma_j^{\text{day}} \varepsilon_{j,t},
\]
and the spot path is constructed through exponential accumulation:
\[
S_{j,t} = S_{j,0}\exp\left(\sum_{u=1}^{t} r_{j,u}\right).
\]

This is useful because the Student-\(t\) innovations generate heavier tails than a Gaussian model, which makes downstream VaR and backtesting behaviour more realistic.

\subsection{Realised Volatility Proxy}

The option-pricing layer uses a realised-volatility proxy. For each factor:
\[
\widehat{\sigma}_{j,t}^{(20)} =
\sqrt{252}\,\mathrm{Std}\left(r_{j,t-19},\dots,r_{j,t}\right),
\]
where the rolling standard deviation is computed over a 20-business-day lookback. This realised volatility is used as the volatility input for Black-Scholes in the listed-option module.

\subsection{Member Archetypes and Positions}

The engine uses ten stylised members, including directional macro, relative value, concentrated volatility, diversified, and weak-liquidity members. The design purpose is not only to create heterogeneity of exposures, but also to ensure that the governance layer has something non-trivial to detect.

At each date \(t\), a member portfolio is represented by quantities \(q_{m,i,t}\). Market value for each position is:
\[
\text{MV}_{m,i,t} = q_{m,i,t} \cdot M_i \cdot S_{\text{und}(i),t},
\]
where \(M_i\) is the contract multiplier and \(\text{und}(i)\) is the underlying risk factor for instrument \(i\).

\subsection{Collateral}

Each member posts a mixture of cash and government bond collateral. Let:
\begin{align*}
C_{m,t}^{\text{cash}} &:= \text{cash collateral},\\
C_{m,t}^{\text{bond}} &:= \text{government bond collateral},\\
h_{m,t} &:= \text{bond haircut}.
\end{align*}

Then posted margin after haircut is:
\[
\text{PM}_{m,t}
= C_{m,t}^{\text{cash}} + C_{m,t}^{\text{bond}}(1-h_{m,t}).
\]

This quantity is what the adequacy layer compares against required initial margin.

\section{Pricing and Daily P\&L}

\subsection{Futures Pricing}

For a futures contract \(i\) on underlying \(j\), the engine uses a simple linear mark:
\[
P_{i,t}^{\text{future}} = M_i S_{j,t}.
\]

For quantity \(q_{m,i,t}\), the contribution to portfolio value is:
\[
V_{m,i,t}^{\text{future}} = q_{m,i,t} P_{i,t}^{\text{future}}.
\]

\subsection{Listed Option Pricing}

For listed options, the engine uses the Black-Scholes model. Let \(S\) be underlying spot, \(K\) the strike, \(T\) time to expiry in years, \(r\) the risk-free rate, and \(\sigma\) the volatility input. Define:
\[
d_1 = \frac{\ln(S/K) + \left(r + \tfrac{1}{2}\sigma^2\right)T}{\sigma\sqrt{T}},
\qquad
d_2 = d_1 - \sigma\sqrt{T}.
\]

Then the call and put prices are:
\begin{align*}
C(S,K,T,r,\sigma) &= S\Phi(d_1) - Ke^{-rT}\Phi(d_2),\\
P(S,K,T,r,\sigma) &= Ke^{-rT}\Phi(-d_2) - S\Phi(-d_1).
\end{align*}

The marked option value in the engine is:
\[
P_{i,t}^{\text{option}} = M_i \times
\begin{cases}
C(S,K,T,r,\sigma), & \text{if } i \text{ is a call},\\
P(S,K,T,r,\sigma), & \text{if } i \text{ is a put}.
\end{cases}
\]

\subsection{Portfolio Value and Daily P\&L}

Portfolio value is the sum over all open instruments:
\[
V_{m,t} = \sum_{i=1}^{I} q_{m,i,t} P_{i,t}.
\]

The one-day P\&L used throughout the engine is the full-revaluation difference:
\[
\Pi_{m,t} = V_{m,t} - V_{m,t-1}.
\]

This is important: the VaR is not built from a linear factor approximation at the member level, but from the realised full-revaluation P\&L series generated by the pricing layer.

\section{Scenario Engine and Baseline Risk}

\subsection{Historical Return Matrix}

Let the daily return vector be:
\[
R_t = (r_{1,t}, r_{2,t}, \dots, r_{J,t})^{\top}.
\]

The historical scenario engine constructs the matrix of the trailing \(500\) business days:
\[
\mathcal{H}_t = \{R_{t-499}, R_{t-498}, \dots, R_t\}.
\]

This forms the basis for historical simulation VaR.

\subsection{Historical Simulation VaR}

For a fixed member \(m\) and date \(t\), define the trailing member-level P\&L sample:
\[
\mathcal{P}_{m,t} = \{\Pi_{m,t-499}, \dots, \Pi_{m,t}\}.
\]

The \(99\%\) Historical Simulation VaR is:
\[
\text{HSVaR}_{m,t}^{99\%}
=
-Q_{0.01}\left(\mathcal{P}_{m,t}\right),
\]
where \(Q_{0.01}\) is the empirical \(1\%\) quantile. Since losses correspond to negative P\&L, the minus sign converts the lower-tail quantile into a positive capital number.

\subsection{Stressed Window Selection}

The engine also computes a stressed VaR using a \(60\)-day stressed sub-window. The stressed period is selected as the contiguous window with the highest realised SPX volatility. If \(r_{\text{SPX},u}\) denotes SPX returns, the stressed-window endpoint can be expressed conceptually as:
\[
u^{\star}
=
\arg\max_{u}
\mathrm{Std}\left(
r_{\text{SPX},u-59}, \dots, r_{\text{SPX},u}
\right).
\]

The stressed set of dates is then:
\[
\mathcal{S}_t = \{u^{\star}-59,\dots,u^{\star}\}.
\]

The stressed VaR is computed by restricting the member P\&L series to this stressed subset:
\[
\text{SVaR}_{m,t}^{99\%}
=
-Q_{0.01}\left(\{\Pi_{m,u}: u \in \mathcal{S}_t\}\right).
\]

If too few stressed observations overlap the available member P\&L history, the implementation falls back to the ordinary historical VaR.

\subsection{Baseline Margin}

The baseline margin is:
\[
B_{m,t} = \max\left(\text{HSVaR}_{m,t}^{99\%}, \text{SVaR}_{m,t}^{99\%}\right).
\]

This is a simple but important anti-procyclicality device. During calm periods, ordinary historical VaR may fall materially. The stressed VaR acts as a floor tied to a high-volatility regime.

\section{Liquidity Add-On}

\subsection{Motivation}

Pure VaR measures mark-to-market loss but not the cost of actually closing out a portfolio. A CCP margin framework should reflect liquidation cost. The engine therefore augments the baseline margin with a liquidity add-on that contains:

\begin{itemize}
    \item half-spread crossing cost,
    \item market impact, and
    \item liquidation-horizon scaling.
\end{itemize}

\subsection{Spread Cost}

For position \(q_{m,i,t}\), underlying spot \(S_{j,t}\), bid-ask spread in basis points \(b_{j,t}\), and multiplier \(M_i\), the spread-crossing component is:
\[
\text{SC}_{m,i,t}
=
|q_{m,i,t}|\,M_i\,S_{j,t}\,\frac{b_{j,t}}{20000}.
\]

The denominator \(20000\) arises because the engine uses half the full bid-ask spread, and basis points must be converted into decimals.

\subsection{Market Impact}

Let \(\text{ADV}_{j,t}\) be the average daily volume of the relevant underlying, and let \(\alpha\) be the impact coefficient. The market-impact estimate is:
\[
\text{MI}_{m,i,t}
=
\alpha
\sqrt{\frac{|q_{m,i,t}|}{\text{ADV}_{j,t}}}
\,S_{j,t}\,|q_{m,i,t}|\,M_i.
\]

In the implementation, \(\alpha = 0.10\).

The square-root form follows the common empirical intuition that execution cost increases sublinearly with size at first, but materially once trade size becomes large relative to market depth.

\subsection{Liquidation-Horizon Scaling}

If \(\ell_j\) is the liquidation horizon in days for factor \(j\), then the raw liquidity cost is scaled by \(\sqrt{\ell_j}\):
\[
\text{LC}_{m,i,t}
=
\left(\text{SC}_{m,i,t} + \text{MI}_{m,i,t}\right)\sqrt{\ell_j}.
\]

The engine uses:
\[
\ell_j =
\begin{cases}
2, & \text{liquid bucket},\\
3, & \text{semi-liquid bucket},\\
5, & \text{illiquid bucket}.
\end{cases}
\]

\subsection{Member Liquidity Multiplier}

Each member type has an additional liquidity multiplier \(\lambda_m^{\text{liq}}\) to reflect expected liquidation difficulty:
\[
\lambda_m^{\text{liq}} \in \{0.7, 0.8, 1.0, 1.3, 1.5\},
\]
with larger values assigned to concentrated-volatility or weak-liquidity members.

The total liquidity add-on is:
\[
\text{LA}_{m,t}
=
\lambda_m^{\text{liq}}
\sum_{i=1}^{I} \text{LC}_{m,i,t}.
\]

\section{Concentration Add-On}

\subsection{ADV Fraction}

Concentration is measured using position size relative to average daily volume:
\[
f_{m,i,t} = \frac{|q_{m,i,t}|}{\text{ADV}_{j,t}}.
\]

\subsection{Piecewise Concentration Rate}

The engine defines a concentration rate function:
\[
c(f) =
\begin{cases}
0, & f < 0.10,\\
0.10, & 0.10 \le f < 0.20,\\
0.25, & f \ge 0.20.
\end{cases}
\]

For each member-day, the implementation takes the maximum concentration rate across instruments:
\[
c_{m,t}^{\max} = \max_i c(f_{m,i,t}).
\]

The concentration add-on is then:
\[
\text{CA}_{m,t} = c_{m,t}^{\max} \cdot B_{m,t}.
\]

This design means concentration does not stack instrument by instrument; rather, the most severe crowding signal determines the add-on rate applied to baseline margin.

\section{Required Margin and Adequacy}

\subsection{Liquidation-Adjusted Loss}

The central quantity in the project is the liquidation-adjusted loss:
\[
\text{LAL}_{m,t}
=
B_{m,t} + \text{LA}_{m,t} + \text{CA}_{m,t}.
\]

The engine sets:
\[
\text{IM}_{m,t} = \text{LAL}_{m,t}.
\]

Equivalently:
\[
\text{IM}_{m,t}
=
\max\left(\text{HSVaR}_{m,t}^{99\%}, \text{SVaR}_{m,t}^{99\%}\right)
 + \text{LA}_{m,t}
 + \text{CA}_{m,t}.
\]

\subsection{Coverage Ratio and Buffer}

Once posted collateral is known, adequacy is measured through:
\[
\text{Coverage}_{m,t} = \frac{\text{PM}_{m,t}}{\text{IM}_{m,t}},
\qquad
\text{Buffer}_{m,t} = \text{PM}_{m,t} - \text{IM}_{m,t}.
\]

\subsection{Traffic-Light Classification}

The traffic-light rule is:
\[
\text{TL}_{m,t}
=
\begin{cases}
\text{green}, & \text{Coverage}_{m,t} \ge 1.10,\\
\text{amber}, & 1.00 \le \text{Coverage}_{m,t} < 1.10,\\
\text{red}, & \text{Coverage}_{m,t} < 1.00.
\end{cases}
\]

This creates a direct operational interpretation:

\begin{itemize}
    \item green means comfortable over-collateralisation,
    \item amber means coverage is technically sufficient but thin,
    \item red means posted collateral is insufficient.
\end{itemize}

\subsection{Margin Call Amount}

Let the shortfall be:
\[
\Delta_{m,t} = \text{IM}_{m,t} - \text{PM}_{m,t}.
\]

Let \(\tau\) be the operational threshold and \(\mu\) the minimum transfer amount. In the implementation:
\[
\tau = 100{,}000, \qquad \mu = 100{,}000.
\]

The margin call function is:
\[
\text{MC}_{m,t}
=
\begin{cases}
0, & \Delta_{m,t} \le \tau,\\
\max(\Delta_{m,t}, \mu), & \Delta_{m,t} > \tau.
\end{cases}
\]

This prevents operational noise from generating very small calls.

\section{Backtesting and Model Monitoring}

\subsection{Exception Definition}

The actual realised loss on date \(t\) is:
\[
\text{AL}_{m,t} = \max(-\Pi_{m,t}, 0).
\]

Backtesting compares this realised loss against prior-day required margin. The exception indicator is:
\[
I_{m,t}
=
\mathbf{1}\left\{
\text{AL}_{m,t} > \text{IM}_{m,t-1}
\right\}.
\]

An exception therefore means the prior day's margin estimate was not sufficient to cover the next day's realised loss.

\subsection{Rolling Exception Count}

For a \(250\)-day rolling window, the engine computes:
\[
E_{m,t}^{(250)}
=
\sum_{u=t-249}^{t} I_{m,u}.
\]

The exception traffic light is:
\[
\text{BT}_{m,t}
=
\begin{cases}
\text{green}, & E_{m,t}^{(250)} < 3,\\
\text{amber}, & 3 \le E_{m,t}^{(250)} < 5,\\
\text{red}, & E_{m,t}^{(250)} \ge 5.
\end{cases}
\]

\subsection{Interpretation}

The monitoring logic is not purely statistical. If exceptions begin to cluster, this suggests a possible regime-specific weakness in the model, for example:

\begin{itemize}
    \item the VaR window may react too slowly,
    \item stressed-window selection may be insufficiently conservative,
    \item liquidity or concentration add-ons may be too weak, or
    \item the options pricing proxy may fail to capture volatility shocks.
\end{itemize}

\section{Data-Quality and Pricing Controls}

\subsection{Stale Data}

For each factor, the engine flags sequences of zero returns. If:
\[
|r_{j,t}| < \varepsilon
\]
for more than two consecutive business days, a stale-data flag is raised. Conceptually, this is a control against unchanged vendor marks rather than a market event.

\subsection{Outlier Returns}

An outlier-return flag is triggered if:
\[
|r_{j,t}| > 0.15.
\]

\subsection{Implausible Volatility}

The realised-volatility input is flagged if:
\[
\widehat{\sigma}_{j,t}^{(20)} < 0.01
\qquad \text{or} \qquad
\widehat{\sigma}_{j,t}^{(20)} > 3.00.
\]

\subsection{Exposure Jumps}

For member-instrument market values, the engine computes day-over-day changes. If a new change exceeds the historical mean plus three standard deviations,
\[
|\Delta \text{MV}_{m,i,t}| > \mu_{m,i}^{\Delta \text{MV}} + 3\sigma_{m,i}^{\Delta \text{MV}},
\]
an exposure-jump flag is raised.

\subsection{Pricing Validity}

Option prices must remain finite and non-negative:
\[
P_{i,t}^{\text{option}} \in [0,\infty), \qquad
P_{i,t}^{\text{option}} \neq \text{NaN}, \qquad
P_{i,t}^{\text{option}} \neq \infty.
\]

These controls do not change the economics of the model, but they are essential for a governed calculation pipeline.

\section{Governance and Escalation}

\subsection{Rule-Based Escalation Framework}

The governance layer turns quantitative states into actions. Define indicator functions:
\begin{align*}
R_{m,t}^{\text{red}} &= \mathbf{1}\{\text{TL}_{m,t} = \text{red}\},\\
R_{m,t}^{\text{call}} &= \mathbf{1}\{\text{MC}_{m,t} > 0\},\\
R_{m,t}^{\text{bt}} &= \mathbf{1}\{E_{m,t}^{(250)} \ge 4\},\\
R_{m,t}^{\text{conc}} &= \mathbf{1}\left\{\frac{\text{CA}_{m,t}}{\text{IM}_{m,t}} > 0.25\right\},\\
R_t^{\text{stale}} &= \mathbf{1}\{\text{critical stale-data condition holds at } t\}.
\end{align*}

The engine's main escalation rules can then be summarised as:

\begin{center}
\begin{tabular}{>{\raggedright\arraybackslash}p{0.23\textwidth} >{\raggedright\arraybackslash}p{0.40\textwidth} >{\raggedright\arraybackslash}p{0.22\textwidth}}
\toprule
Rule & Quantitative condition & Owner \\
\midrule
Single red breach & \(R_{m,t}^{\text{red}}=1\) & Risk Analyst \\
Consecutive red days & two or more consecutive dates with \(R_{m,t}^{\text{red}}=1\) & Senior Risk Officer \\
Methodology review & \(E_{m,t}^{(250)} \ge 4\) & Model Validation \\
Committee watchlist & \(\text{CA}_{m,t}/\text{IM}_{m,t} > 0.25\) & Risk Committee \\
Provisional run & \(R_t^{\text{stale}}=1\) & Market Data Team \\
Operational margin call & \(R_{m,t}^{\text{call}}=1\) & Margin Operations \\
\bottomrule
\end{tabular}
\end{center}

\subsection{Why Governance Matters Mathematically}

Mathematically, the model is a function:
\[
\mathcal{M}: (\text{market data}, \text{positions}, \text{collateral}, \text{parameters})
\mapsto
(\text{margin}, \text{adequacy}, \text{exceptions}, \text{escalations}).
\]

The governance framework defines how to act when the output of \(\mathcal{M}\) breaches certain thresholds. In other words, governance is the decision layer attached to the model.

\section{Reporting Layer}

The reporting module transforms the state of the engine into artifacts intended for daily and committee consumption.

\subsection{Daily Outputs}

For the latest date \(t\), the daily member-adequacy table contains:
\[
(\text{IM}_{m,t}, \text{PM}_{m,t}, \text{Coverage}_{m,t}, \text{TL}_{m,t}, \text{MC}_{m,t}).
\]

The breach register supplements this with:
\[
(\text{breach type}, \text{severity}, \text{owner}, \text{escalation level}, \text{target resolution date}).
\]

\subsection{Weekly and Monthly Outputs}

The weekly report aggregates:
\begin{itemize}
    \item backtesting exceptions,
    \item concentration events,
    \item stale-data incidents, and
    \item escalation counts.
\end{itemize}

The monthly committee pack adds:
\begin{itemize}
    \item time series of total required margin,
    \item weakest members by average coverage,
    \item coverage-ratio distribution statistics,
    \item exception summaries, and
    \item known limitations and recommended actions.
\end{itemize}

\section{End-to-End Daily Pipeline}

At a high level, the daily pipeline can be described as:

\begin{enumerate}
    \item Generate or load market data, instruments, members, positions, and collateral.
    \item Revalue all portfolios and compute member-level daily P\&L.
    \item Compute historical and stressed VaR for each member.
    \item Add liquidity and concentration penalties.
    \item Compute required margin and compare it with posted collateral.
    \item Generate traffic lights and margin calls.
    \item Backtest prior-day margins against realised losses.
    \item Run data-quality and model-control checks.
    \item Generate escalation events.
    \item Produce daily, weekly, and committee-facing reports.
\end{enumerate}

This can be summarised compactly as:
\begin{align*}
(\text{Data})
&\to
(\text{Pricing})
\to
(\text{VaR})
\to
(\text{Add-ons})\\
&\to
(\text{Adequacy})
\to
(\text{Backtesting and Controls})
\to
(\text{Escalation})
\to
(\text{Reporting}).
\end{align*}

\section{Limitations}

The mathematical structure is coherent, but several simplifications remain:

\begin{itemize}
    \item The margin methodology is stylised and omits many CCP production features such as cross-margin offsets, legal account segregation, and default waterfalls.
    \item Option pricing uses a Black-Scholes framework with realised-volatility proxy rather than a full implied-volatility surface.
    \item Rates futures are handled through a linear approximation rather than a richer term-structure model.
    \item Liquidity and concentration costs are intentionally simple and parameterised rather than calibrated from execution data.
    \item Default management and auction mechanics are not modelled.
\end{itemize}

These limitations do not invalidate the project. Rather, they define its role as a teaching and portfolio platform for quantitative risk concepts.

\section{Possible Extensions}

Natural quantitative extensions include:

\begin{itemize}
    \item filtered historical simulation or expected shortfall in place of pure historical VaR,
    \item scenario-dependent implied-volatility surfaces for listed options,
    \item more realistic liquidation curves and impact calibration,
    \item wrong-way collateral risk and eligibility schedules,
    \item member credit-state modelling, and
    \item richer stress scenarios with cross-asset contagion.
\end{itemize}

\section{Appendix: Parameter Summary}

\begin{center}
\begin{tabular}{>{\raggedright\arraybackslash}p{0.40\textwidth} >{\raggedright\arraybackslash}p{0.20\textwidth} >{\raggedright\arraybackslash}p{0.25\textwidth}}
\toprule
Parameter & Value & Role \\
\midrule
VaR confidence level & \(99\%\) & Tail quantile for margin sizing \\
Historical window & \(500\) days & Historical simulation sample \\
Stressed window & \(60\) days & High-volatility sub-window \\
Risk-free rate & \(4.5\%\) & Black-Scholes input \\
Impact coefficient \(\alpha\) & \(0.10\) & Market-impact scaling \\
Green adequacy threshold & \(1.10\) & Comfortable collateral coverage \\
Amber adequacy threshold & \(1.00\) & Minimum full coverage \\
Margin threshold & \$100{,}000 & Operational trigger threshold \\
Minimum transfer amount & \$100{,}000 & Minimum call size \\
Backtesting warning level & \(3\) & Amber exception level \\
Backtesting critical level & \(5\) & Red exception level \\
Stale-data threshold & \(2\) days & Data-quality control threshold \\
Outlier return threshold & \(15\%\) & Return anomaly control threshold \\
\bottomrule
\end{tabular}
\end{center}

\section{Appendix: Mapping Equations to Code}

\begin{center}
\begin{tabular}{>{\raggedright\arraybackslash}p{0.36\textwidth} >{\raggedright\arraybackslash}p{0.50\textwidth}}
\toprule
Concept & Implementation location \\
\midrule
Synthetic market generation & \texttt{src/portfolio.py::generate\_market\_data} \\
Instrument and member generation & \texttt{src/portfolio.py::generate\_instruments}, \texttt{generate\_member\_profiles} \\
Full revaluation P\&L & \texttt{src/pricing.py::compute\_daily\_pnl} \\
Black-Scholes pricing & \texttt{src/pricing.py::\_bs\_price}, \texttt{price\_instrument} \\
Historical and stressed scenarios & \texttt{src/scenarios.py} \\
Historical VaR and stressed VaR & \texttt{src/margin.py::historical\_simulation\_var}, \texttt{stressed\_var} \\
Liquidity add-on & \texttt{src/liquidity.py::compute\_liquidity\_addon} \\
Concentration add-on & \texttt{src/concentration.py::compute\_concentration\_addon} \\
Coverage ratio and margin call & \texttt{src/controls.py} \\
Backtesting exceptions & \texttt{src/backtesting.py::compute\_exceptions} \\
Escalation rules & \texttt{src/escalation.py} \\
Reports and breach register & \texttt{src/reporting.py} \\
\bottomrule
\end{tabular}
\end{center}

\end{document}
